%% kapa: exam=1 round=29 year=2018 period=2 form=A subject=appraisal_law subject_ko=감정평가관계법규 seq=33 number=33
%% source: Q-Net 감정평가사 기출문제 (cst00309) artlSeq=5202808 file=extracted/1008/5202808/2018년 제29회 감정평가사 1차시험 기출문제/2018년 제29회 감정평가사 1차시험 기출문제/2교시 A형.pdf
%% conversion: kordoc 4.24.51 -> tools/gen1.py (auto); review=proofread-sonnet; flags=none
\begin{question}{33}
부동산등기법상 등기부 등에 관한 설명으로 옳지 않은 것은?
\begin{choices}
  \item 등기부는 토지등기부와 건물등기부로 구분한다.
  \item 법원의 명령 또는 촉탁이 있는 경우에는 신청서나 그 밖의 부속서류를 등기소 밖으로 옮길 수 있다.
  \item 등기관이 등기를 마쳤을 때에는 등기부부본자료를 작성하여야 한다.
  \item 폐쇄한 등기기록의 열람 및 등기사항증명서의 발급은 관할 등기소가 아닌 등기소에 청구할 수 없다.
  \item 등기부의 부속서류가 손상ㆍ멸실의 염려가 있을 때에는 대법원장은 그 방지를 위하여 필요한 처분을 명령할 수 있다.
\end{choices}
\end{question}
\answer{4}
